% Generated draft: requires mathematical and source review.
\documentclass[11pt]{article}
\usepackage[margin=25mm]{geometry}
\usepackage{fontspec}
\setmainfont{Latin Modern Roman}
\usepackage{amsmath,amssymb,mathtools,unicode-math}
\setmathfont{Latin Modern Math}
\usepackage{enumitem,needspace,xurl,hyperref}
\setlength{\parindent}{0pt}
\setlength{\parskip}{5pt}
\setlength{\emergencystretch}{3em}
\newcommand{\N}{\mathbb{N}}
\newcommand{\Z}{\mathbb{Z}}
\newcommand{\Q}{\mathbb{Q}}
\newcommand{\R}{\mathbb{R}}
\newcommand{\C}{\mathbb{C}}
\newcommand{\F}{\mathbb{F}}
\newcommand{\A}{\mathbb{A}}
\newcommand{\PP}{\mathbb{P}}
\newcommand{\E}{\mathbb{E}}
\newcommand{\eps}{\varepsilon}
\newcommand{\dd}{\,\mathrm{d}}
\begin{document}
\section*{UnsolvedMath ID: 700703}
\textbf{Source problem code:} catalogue entry
\subsection{Definitions and mathematical statement}
Let \(F(a,x)\) be the free group on generators \(a,x\), and let a group word \(w(a,x)\) mean a finite word in \(a^{\pm1},x^{\pm1}\).

\textbf{Part (a), attributed in the source to R. C. Lyndon.} The source asks for an algorithm deciding whether the distinguished generator \(a\) is trivial in
\[
G_{n,w}=\langle a,x\mid a^n=1,\ w(a,x)=1\rangle.
\]
The exponent \(n\) is not quantified in the supplied statement, and its domain and role as an algorithmic input are unspecified. Consequently, the following are conditional formulations, rather than an unqualified resolution of the intended question. Under the interpretation \(n\in\Z_{>0}\), define
\[
D_n=\{w\in F(a,x):a=1\text{ in }G_{n,w}\}
=\{w\in F(a,x):a\in\langle\!\langle a^n,w\rangle\!\rangle_{F(a,x)}\},
\]
where \(\langle\!\langle a^n,w\rangle\!\rangle_{F(a,x)}\) is the smallest normal subgroup containing \(a^n\) and \(w\). A fixed-parameter reading asks whether, for every fixed positive integer \(n\), there exists an algorithm that halts on every input word \(w\) and decides membership in \(D_n\). A uniform reading asks whether one algorithm halts on every input pair \((n,w)\), with \(n\in\Z_{>0}\), and decides whether \(w\in D_n\). The supplied record does not distinguish these readings or exclude an omitted restriction on \(n\).

\textbf{Part (b), attributed in the source to R. C. Lyndon and marked as archived solved.} Using the conventions
\[
u^v=v^{-1}uv,\qquad [u,v]=u^{-1}v^{-1}uv,
\]
let
\[
G=\langle a,x\mid a^5=1,\ a^{x^2}=[a,a^x]\rangle.
\]
The question is whether the image of \(a\) in \(G\) is different from the identity. The conjugation and commutator conventions are made explicit here; they are not explicitly defined in the supplied record and require confirmation against the source. The supplied archive reports an affirmative answer, attributed to D. N. Azarov; this is a source-reported solved label, not an independently verified conclusion of this draft.
\subsection{Short English statement}
Part (a) asks whether an algorithm can decide, from the additional relator \(w(a,x)\), whether \(a\) becomes trivial in the presented group. The source does not specify the domain or quantification of \(n\), so the precise algorithmic target remains unresolved. Part (b) asks whether \(a\) is nontrivial in a particular presented group and is reported as solved affirmatively in the supplied archive.
\subsection{Sources}
\begin{itemize}
\item Supplied local UnsolvedMath registry record: permanent ID \(700703\), exact source code \texttt{catalogue entry}, title \emph{Kourovka Notebook Problem 11.10}. The registry attributes the proposal to V. V. Bludov, with proposed year 1990; the statement attributes both parts to R. C. Lyndon. Their respective roles are retained as supplied.
\item \emph{The Kourovka Notebook (Unsolved Problems in Group Theory)}, 11th issue (1990), Problem 11.10; supplied background identifies an update in the 21st issue (2026), \texttt{arXiv:1401.0300v48}. The supplied collection description names E. I. Khukhro and V. D. Mazurov as editors. These references have not been independently consulted.
\item Supplied Archive of Solved Problems entry for part (b): D. N. Azarov, in \emph{Algebraicheskije Systemy}, Ivanovo, 1991, pp.\ 4--5 (Russian). The archive reports that the assignments
\[
a\longmapsto(1\ 3\ 5\ 2\ 4),\qquad x\longmapsto(1\ 2\ 4\ 3\ 5)
\]
extend to a homomorphism from \(G\) onto the alternating group \(\mathbb{A}_5\), and reports an affirmative answer to part (b). This archive assertion is preserved as supplied, without independent verification.
\item The supplied registry status is \texttt{partially\_solved}. Its status review is dated \texttt{2026-10-16}, lists evidence as \texttt{published}, and states that solved parts occur in the archive while remaining parts occur in the main section. This assessment is supplied dated curation, not independently verified mathematical or literature review.
\end{itemize}
\end{document}
